\section{Two separate norms and a balanced Pell quotient}
\label{sec:low-degree}

The direct positive-variable route admits a substantial reduction that is
independent of the nine-variable packing construction.  The recommended
polynomial has twelve positive unknowns, exact degree $242$, and $94$ binary
operations.  A cheaper companion has degree $258$ and $89$ operations.
The incoming wave-7 comparison is twelve positive unknowns, degree $897$,
and $108$ operations.  A further thirteen-positive companion has degree
$132$ and $83$ operations.  The twelve-positive degree-$314$, $91$-operation
intermediate is retained because its reconstruction makes the mechanism
especially transparent.

The new mechanism keeps two quadratic norms separate.  One combines the
main Pell root and the first Pell root while absorbing the index quotient.
The other removes the strong Pell root.  A known residue of the auxiliary
Pell ordinate then permits an exact polynomial factor to be removed, with
its size chosen to balance the two remaining terms.  All coefficients are
the complete compact tuple in the coefficient appendix, with the already
shifted high mask.  There are no additional existential coefficient ports.

\subsection{The literal twelve-variable polynomial}

Supply exactly
\[
 J,F,\alpha,z,i,T,s,w,\eta,\zeta,y,t\in\mathbb Z_{>0}.
\]
Every other letter in this section denotes a polynomial abbreviation.  Set
\begin{align*}
 q&=\beta J+1,&X&=wq,&Y&=sq^3,&E&=XY,\\
 k&=\eta+\zeta,&c&=kY+\eta,&a&=E+Y,&H&=4a+3,\\
 \Delta&=a^2+H,&\ell&=EkY,&A&=1+\ell(\ell+k),&G&=1+\Delta c^2,\\
 b_0&=X+ac,&U&=4q^4,&r&=k-R-1,&S&=UH.
\end{align*}
Here the transport expressions are
\begin{align*}
 R&=((q-1)(q-F)+\alpha)(q^2-1)+(M_C+qM_F)J,\\
 N_t&=(K_0+w)(q-F-\alpha)+q-F-z(q-1).
\end{align*}
Define the first norm by
\begin{align}
 W&=Ub_0+H\bigl(U(t+\ell)+r\bigr),\notag\\
 \mathcal N&=(W^2+U^2G-S^2A)^2-4W^2U^2G.
 \label{eq:two-root-norm}
\end{align}
For the second norm put
\begin{align*}
 I&=ic^2,&D&=1+\Delta I^2,&Q&=(\Delta I)^2,&d&=Q-1,\\
 v&=cT,&b&=c+RD.
\end{align*}
Choose one of the following two pairs of polynomial abbreviations:
\[
 (M,C)=\bigl((\Delta ic)^2,c^2\bigr)
 \quad\hbox{or}\quad
 (M,C)=(Q,1).
\]
In both cases $Q=MC$.  Finally set
\begin{align}
 Z&=C(Dv^2+b^2-1)+dy(2-My),\notag\\
 \mathcal J_M&=Z^2-4C^2Dv^2b^2,\notag\\
 \boxed{P_M=N_t(1+\mathcal N)(1+\mathcal J_M)-1=0.}
 \label{eq:twelve-balanced}
\end{align}
The first choice is the degree-$242$ equation; the second is the
$258$-degree companion.  Both are ordinary integer polynomials.  The
balanced evaluator computes $M=(\Delta I)(\Delta i)$ and $Q=Mc^2$,
which is the same polynomial definition with one multiplication saved.

\begin{theorem}
For each of the two fixed polynomials in \eqref{eq:twelve-balanced}, ZFC
proves that it has no solution in twelve strictly positive integers if and
only if $\operatorname{Con}(\mathrm{ZFC})$.  Their exact degree/operation
pairs are $(242,94)$ and $(258,89)$, respectively.
\end{theorem}

The theorem composes the reduction below with the inherited native-window
compiler and K215 syntactic-refutation calibration.  The complete new
argument is supplied here; the source texts for the signed Pell recovery,
canonical compiler witnesses, and auxiliary congruences are retained in
\texttt{low\_degree/upstream/}, with exact pins.  This is an ordinary
mathematical proof formalizable in ZFC, rather than a new exported
object-language derivation or a newly checked end-to-end proof-assistant
artifact.

\subsection{Units, positivity, and recovery of the strong root}

Both $\mathcal N$ and $\mathcal J_M$ are squares modulo $4$.  If
\eqref{eq:twelve-balanced} holds, its three integer factors are units.
The value $-1$ for either $1+\mathcal N$ or $1+\mathcal J_M$ would force
the corresponding norm to equal $-2$, impossible modulo $4$.  Thus
\begin{equation}
 N_t=1,\qquad \mathcal N=\mathcal J_M=0.
 \label{eq:low-units}
\end{equation}
This obtains the transport interface before any Pell index argument.  Its
inherited bounds are
\begin{equation}
\begin{gathered}
 q\ge16,\quad C_w=q-F-\alpha>0,\quad 3q<R<q^4-2,\\
 X\ge q,\quad Y\ge q^3,\quad E\ge q^4>R+2,\qquad
 k\ge2,\quad kY<c<k(Y+1).
\end{gathered}
\label{eq:low-outer}
\end{equation}
For example, $C_w\le0$ would imply
$N_t\le q-F-z(q-1)=1-F-(z-1)(q-1)\le0$.  The remaining bounds are the
authentic whole-tuple mask bounds.  Consequently $M,C,v,b$ are positive,
$Q>1$, and $M\ge2$.  All denominators used next are nonzero.

From $\mathcal J_M=0$ define the rational number
\[
 f=\frac{Z}{2Cvb}.
\]
It satisfies $f^2=D$, so it is an integer.  The sign of this recovered
integer is not prescribed.  Put $y_{\rm old}=My-1>0$.  The exact polynomial
identity
\begin{equation}
 MZ=Q(Dv^2+b^2)-1-d(My-1)^2
 \label{eq:quotient-identity}
\end{equation}
then gives
\begin{equation}
 f^2-\Delta I^2=1,\qquad
 d y_{\rm old}^2=Q(vf-b)^2-1,
 \qquad vf-b=c(Tf-1)-Rf^2.
 \label{eq:low-strong-aux}
\end{equation}
These are the normalized strong and auxiliary Pell equations.  The signed
auxiliary theorem uses $|f|$ as its modulus and covers a negative $f$.
To match the preserved ordinary-strong statement explicitly, take
$i_{\rm ordinary}=\Delta i>0$; then
$(i_{\rm ordinary}c^2)^2=\Delta(f^2-1)=Q$.  No positive sign for the
reconstructed $f$ is silently assumed.

\subsection{A two-root carrier restores the main square}

Write $g_0=\sqrt G>0$ and $\tau_0=\sqrt A>0$.  There is the exact real
factorization
\begin{equation}
 \mathcal N=\prod_{\sigma,\epsilon\in\{-1,1\}}
       (W-\sigma Ug_0-\epsilon S\tau_0).
 \label{eq:main-four-factors}
\end{equation}
Since $r=k-R-1\ge1-R$, $U>R$, and $t\ge1$, we have $t+r/U>0$ and
\[
 W=Ub_0+S(\ell+t+r/U)>0.
\]
A vanishing factor supplies real roots $g^2=G$, $\tau^2=A$ with
$W=Ug+S\tau$.  Squaring isolates
\[
 2WU g=W^2+U^2G-S^2A.
\]
Thus $g$ is rational, hence integral; then $\tau$ is integral as well.
This argument remains valid when the two radicands have the same square
class.

The sign $g<0$ is impossible.  Indeed
\[
 \tau_0<\ell+k,\qquad
 \frac{b_0+g_0}{H}>\frac{2ac}{4a+3}>\frac c3>k,
\]
since $a>1$ and $Y>3$.  For $g=-g_0$ the factor equation would give
\[
 \tau=\ell+t+r/U+(b_0+g_0)/H>\ell+k>\tau_0,
\]
a contradiction.  Therefore $g=g_0>0$.  Reducing the carrier identity
modulo $U$ proves
\begin{equation}
 U\mid H(k-R-1).
 \label{eq:low-weak-divisibility}
\end{equation}
We have not yet canceled $H$ or recovered the omitted index quotient.

\subsection{Recover the exact index before the deleted quotients}

Put $A_0=a+2$ and $P_0=2EY+1$.  Standard Pell classification gives
positive indices $p,m,n$ with
\begin{align*}
 g&=\chi_{A_0}(p),&c&=\psi_{A_0}(p),\\
 |f|&=\chi_{A_0}(m),&I&=\psi_{A_0}(m),\\
 |\tau|&=\chi_{P_0}(n),&k&=2\psi_{P_0}(n).
\end{align*}
The first-Pell discriminant is $EY(EY+1)$, with fundamental
integer-coordinate unit $(P_0,2)$.  Because $P_0>A_0$ and $c>k$,
monotonicity gives $n<p$.  Also $c>1$ implies $p\ge2$, whence
$c\ge2A_0>2R$, and Pell growth gives $c>2p$.

Strong divisibility and $c^2\mid I$ give $p\mid m$.  Write $m=pj$.
Then
\[
 I=c\psi_g(j),\qquad \psi_g(j)\equiv jg^{j-1}\pmod c.
\]
Since $g^2\equiv1\pmod c$, $g$ is invertible modulo $c$, so $c\mid j$.
Consequently $c\mid m$ and $m\ge pc>2p$.

For clarity, the signed auxiliary-index recovery used at this point does
not depend on either deleted quotient.  Put $M_0=\Delta I$ and
$V=c(Tf-1)-Rf^2$.  The second equation in \eqref{eq:low-strong-aux} is
\[
 M_0^2V^2-(M_0^2-1)y_{\rm old}^2=1.
\]
Its Pell index $j_0$ is odd because its abscissa is divisible by $M_0$.
For the integer odd-index polynomials $Q_j$,
\begin{align*}
 \chi_{M_0}(2j+1)&=M_0Q_j(M_0^2),\\
 Q_j(1-A_0^2)&=(-1)^j\psi_{A_0}(2j+1),\\
 Q_j(0)&=(-1)^j(2j+1).
\end{align*}
Modulo $|f|$, $M_0^2\equiv1-A_0^2$ and $V\equiv-c$.  Squaring and
doubling, and reducing the signed abscissa index using $0<2p<m$, yields
$j_0\equiv\pm p\pmod m$, hence modulo $c$.  Modulo $c$,
$M_0^2\equiv0$, $f^2\equiv1$, and $V\equiv-R$, so
$j_0\equiv\pm R\pmod c$.  Since $0<p,R<c/2$, it follows that
\begin{equation}
 p=R.
 \label{eq:low-index}
\end{equation}
This is the signed index argument reproduced in the preserved sources;
its inputs have all been verified before using $h$ or $\gamma$.

Let $\mu=U/\gcd(U,H)$.  Since $H$ is odd and $H\equiv3\pmod q$, the
only possible common prime of $q$ and $H$ is $3$.  If $3\mid q$, then
$27\mid Y$ and $H\equiv3\pmod{27}$, so the $3$-adic valuation is one.
Therefore
\[
 \gcd(U,H)\le3,\qquad \mu\ge4q^4/3>R.
\]
The first Pell recurrence gives $k\equiv2n\pmod{4EY}$.  The modulus
$4EY=4ws^2q^7$ is a multiple of $U$.  Thus
\eqref{eq:low-weak-divisibility} gives $2n\equiv R+1\pmod\mu$.
But $1\le n<R$ implies $|2n-R-1|<R<\mu$.  Hence
\[
 2n=R+1,\qquad r=k-R-1\equiv0\pmod{4EY}.
\]
In particular $h=r/E$ and $r/U$ are integers.  Moreover $n\ge25$ and
$\psi_{P_0}(n)>n$ make $h>0$.  Dividing the carrier by $U$ gives
\[
 g=b_0+H(t+\ell+r/U-\tau),
\]
so $\gamma=(g-b_0)/H$ is integral.  Since $n\ge2$ gives
$k\ge4P_0>X$, we have $c>X$; since $\Delta>(a+1)^2$,
$g-b_0>c-X>0$.  Thus $\gamma>0$.

All normalized compiler equations have now been recovered, including
$k-hE-R=1$ and the positive main quotient.  Their signed auxiliary
soundness theorem gives a blank K215 halt, hence a finite syntactic ZFC
refutation.  The proof does not require making the reconstructed $f$
positive on the same assignment.

\subsection{Canonical completeness and the removed Pell residue}

Start with a canonical positive normalized compiler witness.  It has
$p=R=2n-1$, $h>0$, and $\gamma>0$.  Its first Pell recurrence gives
$U\mid r$.  Define
\begin{equation}
 t=\gamma+\tau-\ell-r/U.
 \label{eq:low-t-complete}
\end{equation}
This is a positive integer: since $\tau^2=\ell^2+\ell k+1$ and
$\ell>k$,
\[
 \tau-\ell=\frac{\ell k+1}{\tau+\ell}
 >\frac{\ell k}{2\ell+k}>\frac k3>\frac rU.
\]
The carrier identity follows, so $\mathcal N=0$.

For the new ordinate quotient, one may replace the canonical auxiliary
witnesses explicitly.  The authentic population decoding gives
$R\equiv3\pmod4$.  To identify the precise inherited fact, the native
radix is divisible by $4$, $q=B^N$, $J\equiv1\pmod4$, and
$M_C\equiv2\pmod4$.  The low mask forces $C_w\equiv1\pmod4$; the
formula for $R$ gives $R\equiv C_w+2\equiv3\pmod4$.  This argument is
preserved in the wave-7 squared-modulus source and applies to the whole
compact tuple here.

Keep the outer coordinates and the original positive $f,i$.  Put
$m_0=\Delta I$, so $Q=m_0^2$, and choose
\begin{equation}
 V=\chi_{m_0}(R)/m_0,\qquad
 y_{\rm old}=\psi_{m_0}(R).
 \label{eq:canonical-auxiliary}
\end{equation}
The odd-index polynomial makes $V$ a positive integer.  Its identities give
\begin{equation}
 V\equiv-c\pmod{f^2},\qquad
 V\equiv-R\pmod c,
 \qquad y_{\rm old}\equiv-1\pmod Q.
 \label{eq:canonical-residues}
\end{equation}
For the first use $m_0^2\equiv1-A_0^2\pmod{f^2}$ and
$(R-1)/2$ odd; for the second use $c\mid m_0$ and the constant term
of $Q_{(R-1)/2}$.  For the third, $\psi_x(R)$ is an even polynomial
in $x$, with constant term $(-1)^{(R-1)/2}=-1$.  The recurrence
$\psi_x(j+1)=2x\psi_x(j)-\psi_x(j-1)$ proves this assertion by induction.

Because $\gcd(c,f)=1$ and $f^2\equiv1\pmod c$, the numerator below is
divisible by $cf$ (in fact by $cf^2$).  Thus
\[
 T=\frac{V+c+Rf^2}{cf}>0
\]
is an integer, and $V=c(Tf-1)-Rf^2$.  The Pell identity gives the auxiliary
equation.  This replaces only $T$ and the auxiliary ordinate; it leaves
\eqref{eq:low-t-complete} unchanged.  No assumption of cofinal auxiliary
witnesses is needed.

For either selected $M\mid Q$, define
\[
 y=(y_{\rm old}+1)/M\in\mathbb Z_{>0}.
\]
Identity \eqref{eq:quotient-identity} gives $\mathcal J_M=0$.  All twelve
supplied coordinates are positive, proving completeness.  Negating the
existence equivalence proves the stated consistency equivalence in ZFC.

\subsection{Exact degrees and the intermediate companions}

Write the leading homogeneous forms as
\begin{gather*}
 q_0=\beta J,\quad X_0=wq_0,\quad Y_0=sq_0^3,\quad E_0=X_0Y_0,\quad
 k_0=\eta+\zeta,\quad c_0=k_0Y_0,\quad \ell_0=E_0c_0,\\
 H_0=4E_0,\quad U_0=4q_0^4,\quad I_0=ic_0^2,\quad
 R_0=q_0^3(q_0-F),\quad (N_t)_0=w(q_0-F-\alpha)-zq_0.
\end{gather*}
The large apparent degree of \eqref{eq:two-root-norm} cancels by an exact
polynomial identity.  Set
\begin{align*}
 D_0&=b_0^2-G=X^2+2Xac-Hc^2-1,&Z_0&=Ut+r,\\
 C_2&=U^2D_0+2UHb_0Z_0
       +H^2(2U\ell Z_0+Z_0^2-U^2\ell k-U^2).
\end{align*}
Then
\begin{equation}
\begin{split}
 \mathcal N={}&4U^4H^2\ell^2D_0+4U^2Hb_0\ell C_2+C_2^2\\
 &\hspace{1cm}-4U^4H^2G(\ell k+1).
\end{split}
\label{eq:norm-degree-identity}
\end{equation}
Here $\deg D_0\le16$ and $\deg C_2\le32$.  The four summands have
respective degree bounds $66,68,64,62$.  Only the second reaches $68$,
with nonzero highest part
\begin{equation}
 \mathcal N_{68}=2^{16}w^6s^9k_0^3q_0^{49}(2t-k_0).
 \label{eq:low-norm-leader}
\end{equation}
This proves the upper and lower degree statements by polynomial identities,
without treating a sampled diagonal degree as an upper bound.

For the full-modulus choice $M=Q$, the two summands defining $Z$ have
maximum degrees $76$ and $94$.  The latter has highest part
$-Q_0^2y^2$, where $Q_0=E_0^4I_0^2$.  Thus $\mathcal J_M$ has degree
$188$, and $P_M$ has degree $68+188+2=258$, with highest part
\[
 \mathcal N_{68}Q_0^4y^4(N_t)_0\ne0.
\]
For the balanced choice, $\deg M=36$ and $\deg C=10$.  The first summand
of $Z$ has degree $86$ and highest part $c_0^2R_0^2E_0^4I_0^4$; the
second has degree $84$.  The subtracted term of $\mathcal J_M$ has degree
$142$.  Hence $\deg\mathcal J_M=172$, and the highest part of $P_M$ is
\[
 \mathcal N_{68}c_0^4R_0^4E_0^8I_0^8(N_t)_0\ne0,
\]
of degree $242$.

The degree-$314$ intermediate uses the original ordinate directly:
\begin{align*}
 e&=Q(Dv^2+b^2)-1,& B_{\mathrm{aux}}&=2Qvb,\\
 \mathcal J&=(e-dy^2)^2-B_{\mathrm{aux}}^2D,&
 P_{314}&=N_t(1+\mathcal N)(1+\mathcal J)-1.
\end{align*}
It reconstructs $f=(e-dy^2)/B_{\mathrm{aux}}$, with no canonical residue
restriction.  Its inner norm has degree $244$ and highest part
$E_0^{16}I_0^{12}R_0^4$, proving total degree $314$.

Finally, retaining $f>0$ gives the thirteen-positive companion
\begin{align*}
 N_s&=f^2-\Delta I^2,&V&=c(Tf-1)-Rf^2,\\
 N_a&=\Delta(f^2-1)(V^2-y^2)+y^2,&
 P_{132}&=N_t(1+\mathcal N)N_sN_a-1.
\end{align*}
A zero forces all factors to be integer units and $\mathcal N=0$.
The elementary negative-Pell obstruction
\[
 x^2-(m^2-1)y^2\ne-1\qquad(m\ge2)
\]
first rules out $N_s=-1$.  Then $\Delta(f^2-1)=(\Delta I)^2$ is a
square, so the same obstruction rules out $N_a=-1$.  Thus $N_t=N_s=N_a=1$.
For the obstruction, a hypothetical nonnegative solution has
$(m-1)y<x<my$; multiplying by $m-\sqrt{m^2-1}$ gives a smaller positive
ordinate $my-x$, contradicting descent.  The remaining proof is the
same carrier/index reconstruction above.  Degrees of the factors are
$2,68,34,28$, and their nonzero leading product is
\[
 -\mathcal N_{68}E_0^4I_0^2c_0^2T^2f^4(N_t)_0,
\]
so the exact degree is $132$.

The literal DAGs have $94$, $89$, $91$, and $83$ live binary gates for
$P_{242},P_{258},P_{314},P_{132}$, respectively.  Independent finite checks
include full polynomial arrays over three primes using the actual fixed
coefficients, four-factor multiplication in the formal quadratic algebra,
unit-guard congruences, the exact quotient identity, and nontrivial canonical
Pell reconstructions.  These validate the delivered arithmetic; the universal
proof and the identified inherited compiler arguments establish the theorem.
